\section{The Graph That Does Not Compose}
\label{sec:ch07-the-graph-that-does-not-compose}

\index{composition!what the composer reads|(}%
\index{wgc@\code{wgc}!installing it}%
The composer is a command-line tool called \code{wgc}, and it comes from npm
rather than NuGet, so this is the first thing in the book that needs Node
installed:

\begin{minted}{text}
npm install -g wgc
\end{minted}

The version table in appendix~\ref{app:toolchain} pins the release this book's
error text was captured with, and a later one may word an error differently,
so install that one rather than whatever is current.

\index{repository layout}%
It takes one input file, and where that file lives is the first time the
layout of the whole project matters. Everything this book builds sits in one
tree: one folder per service under \code{src}, so \code{src/Sessions} is the
project chapter~\ref{ch:hot-chocolate-zero} created and the next two services
go beside it; \code{graph} for what the composer reads, which is this file
and nothing generated; and \code{router} for the one file
chapter~\ref{ch:enter-the-router} adds. Here is the whole of the input, for a
graph with one subgraph in it:

\begin{minted}{yaml}
version: 1
subgraphs:
  - name: sessions
    routing_url: http://localhost:5001/graphql
    schema:
      file: ../src/Sessions/schema.graphql
\end{minted}

\index{wgc@\code{wgc}!reads files, not services}%
Read what that file names. A schema on disk, and a url the composer does not
call. \code{routing\_url} is not where the composer goes to find your schema;
it is what the composer writes into its output so that a router knows where to
send a request later. Nothing here talks to a running service, and your service
does not have to be running. That is the first fact about composition and most
of what follows depends on it.

\index{wgc@\code{wgc}!the compose command}%
Then the command, from the directory holding that file:

\begin{minted}{text}
wgc router compose -i graph.yaml -o router.json
\end{minted}

\index{composition!the first run fails}%
One subgraph. The one this book has been building since
chapter~\ref{ch:hot-chocolate-zero}, exported by the service itself, with the
federation additions chapter~\ref{ch:first-subgraph} put in and nothing else
changed. It does not compose:

\begin{minted}[fontsize=\footnotesize,breakanywhere]{text}
The subgraph "sessions" could not be federated for the following reasons:
The 1st instance of the directive "@cost" declared on coordinates "Query.node" is invalid for the following reason:
 The value ""10"" provided to argument "@cost(weight: ...)" is not a valid "Int!" type.
The 1st instance of the directive "@cost" declared on coordinates "Query.nodes" is invalid for the following reason:
 The value ""10"" provided to argument "@cost(weight: ...)" is not a valid "Int!" type.
The 1st instance of the directive "@listSize" declared on coordinates "Query.sessions" is invalid for the following
reason:
 The definition for "@listSize" does not define the following argument that is provided: "slicingArgumentDefaultValue…
The 1st instance of the directive "@cost" declared on coordinates "Query.sessions" is invalid for the following
reason:
 The value ""10"" provided to argument "@cost(weight: ...)" is not a valid "Int!" type.
The 1st instance of the directive "@cost" declared on coordinates "Query.sessionById" is invalid for the following
reason:
 The value ""10"" provided to argument "@cost(weight: ...)" is not a valid "Int!" type.
The 1st instance of the directive "@cost" declared on coordinates "Mutation.rescheduleSession" is invalid for the
following reason:
 The value ""10"" provided to argument "@cost(weight: ...)" is not a valid "Int!" type.
The 1st instance of the directive "@cost" declared on coordinates "Session.speaker" is invalid for the following
reason:
 The value ""10"" provided to argument "@cost(weight: ...)" is not a valid "Int!" type.
\end{minted}

\index{composition!errors name a coordinate}%
Three things about that block before anything about why. The doubled quotes in
\code{""10""} are the composer quoting a value that is itself a quoted string,
not a typo. The line about \code{@listSize} stops mid-word because the tool
draws a fixed-width box around its errors and truncates to fit, so the argument
name really does arrive cut off. And every error names a schema coordinate
rather than a line number, which is the right choice for a document nobody
wrote by hand.

\index{composition!nothing here is about federation}%
Now read the vocabulary. No key, no entity, no seam, no subgraph but your own.
Six rejected cost weights and one rejected argument, seven refusals over one
service, and not one of them is about the thing this part of the book is for. I
had expected the first composition of this graph to fail. I had not expected it
to fail without mentioning federation.

\subsection{Nobody Wrote Either Directive}

\index{cost@\code{"@cost}!nobody configured it}%
Go back and look at where \code{@cost} came from. It has been in this schema
since chapter~\ref{ch:hot-chocolate-zero}, where it appeared on both fields
\code{Query} had at the time and was named as something nobody asked for. Hot
Chocolate applies a default cost weight with no configuration at all, to the
fields that have a resolver behind them. Six of them are in this schema now,
and the six are exactly the six the composer rejected.

\index{listSize@\code{"@listSize}!comes from the paging convention}%
\code{@listSize} arrived with paging in
chapter~\ref{ch:data-without-n-plus-1}. It carries the two numbers that
chapter's \code{[UsePaging]} attribute set, and it carries them because the
same cost analyzer wants to know how big a page can get.

\index{composition!rejects what the library generated}%
So the schema that fails to compose contains, in the places that fail, nothing
you typed. Chapter~\ref{ch:first-subgraph} ended on the observation that the
package writes directives into your public API without asking and that you
should read the exported file after switching federation on. This is the same
lesson arriving with a bill attached, and from a different library.
\index{composition!what the composer reads|)}%
